El artículo \emph{The value of our personal data in the Big Data and the Internet of all Things Era}, escrito por Anahiby Becerril, aborda cómo la información personal se ha convertido en un recurso altamente mercantilizado dentro de la economía digital contemporánea.
Los individuos han dejado de ser simples usuarios de la tecnología para transformarse en activos intangibles que entregan continuamente fragmentos de su vida a cambio del acceso a plataformas en línea.
Considero que revisar este texto es de una importancia crítica en la actualidad, debido a que desde su publicación original en 2018, el problema descrito por la autora no ha hecho más que agudizarse y profundizarse en todos los niveles de nuestra interacción tecnológica.
Si ya en 2018 era un problema grave y resultaba urgente que la gente tomara conciencia y se empoderara frente a las corporaciones, hoy en pleno 2026, buscar esta toma de conciencia es más apremiante que nunca.
Durante estos últimos años, la producción de datos ha seguido aumentando exponencialmente y hemos sido testigos de las devastadoras consecuencias de esta hiperconectividad.
Hemos visto el impacto de entidades como \emph{Cambridge Analytica} en la manipulación e influenciamiento de poblaciones enteras (\cite{cadwalladr2018cambridge}), y más aún, estamos presenciando el surgimiento de las grandes empresas de Inteligencia Artificial, las cuales entrenan sus modelos masivos con literalmente toda la información digital disponible, frecuentemente sin transparencia ni compensación justa (\cite{grynbaum2023nyt}, \cite{genai}, \cite{guild}).

En su obra, Becerril expone que cada interacción diaria con dispositivos móviles, sensores o redes sociales deja un rastro digital que las empresas recolectan y procesan para obtener enormes ventajas competitivas y ganancias millonarias.
Esta dinámica de la cuarta revolución industrial se basa en predecir comportamientos y perfilar a los individuos, quienes ignoran en gran medida el valor monetario real de su información, como su ubicación o hábitos financieros.
Estas ideas del artículo se ven perfectamente materializadas y amplificadas por eventos recientes que ilustran los peligros de esta asimetría de poder.
El infame escándalo de \emph{Cambridge Analytica} evidenció cómo se extrajeron indebidamente los datos privados y perfiles psicológicos de hasta 87 millones de usuarios de Facebook para crear mensajes ultrapersonalizados que influyeron en campañas políticas y elecciones (\cite{cadwalladr2018cambridge}).
Hoy, este modelo extractivista ha sido heredado y escalado por compañías como Anthropic y OpenAI, las cuales consumen todo tipo de fuentes de información, incluyendo repositorios de código abierto masivos y plataformas de internet, para nutrir y entrenar a sus grandes modelos de lenguaje (LLMs) (\cite{grynbaum2023nyt}).

Como ya mencioné, suscribo totalmente los planteamientos del artículo, ya que ejemplos como los ya mencionados demuestran indiscutiblemente que los datos tienen un valor intrínseco colosal que las grandes corporaciones explotan de forma unilateral para su propio beneficio monetario o de poder.
Sin embargo, me parece que la percepción de la gente no ha cambiado mucho y la mayoría de los usuarios sigue cediendo su información en todo momento sin cuestionarlo.
Un ejemplo alarmante de esto es el uso sistemático de las conversaciones de los usuarios con los LLMs para refinar, afinar y entrenar las próximas iteraciones de estos modelos.
Un caso reciente y paradigmático en este 2026 es el escándalo que involucra a OpenAI y la resolución del histórico problema matemático de Navier-Stokes.
Se reportó que los agentes de IA de OpenAI lograron una demostración para este problema luego de que el modelo Codex presuntamente asimilara la información, borradores y sesiones de trabajo de dos matemáticos (uno de ellos investigador de Anthropic) que estaban buscando su propia prueba (\cite{quanta2026navier}).
Todo esto aparentemente ocurrió sin el consentimiento explícito de los matemáticos sobre el uso de sus datos específicos de investigación, una situación que en otros contextos académicos y profesionales se consideraría simplemente plagio.

El objetivo principal que la autora intenta plantear es convencernos de que, como sociedad, debemos darnos cuenta de la importancia de los datos que producimos.
Intenta persuadirnos de buscar mecanismos tecnológicos y legales para empoderarnos y lograr un control real de los mismos, evitando que se sigan cediendo por ignorancia u omisión a los gigantes corporativos.
Todo este panorama se relaciona enormemente con la materia de Fundamentos de Bases de Datos.
En principio, esta materia nos prepara profesionalmente para estructurar, gestionar y manejar grandes volúmenes de información, potencialmente provenientes de usuarios que, de manera consciente o no, nos encomiendan el resguardo de su vida digital.
En cierta forma, este conocimiento técnico nos hace directamente responsables de asegurar un uso correcto, seguro y ético de dichos datos.
Fomentar una mala práctica, una pobre seguridad o un mal diseño de base de datos puede contribuir activamente a agravar el problema de indefensión ciudadana descrito por la autora.
Por ello, mientras la temática central del texto es la mercantilización y el valor económico de nuestra información, los temas laterales como la privacidad, el Internet de las Cosas y la gobernanza global de los datos impactan directamente nuestra práctica profesional y la forma en que concebimos la arquitectura de la información.

El artículo me aporta información certera y precisa sobre un problema del que soy consciente, aunque admito que no siempre soy consecuente en la práctica.
En mi vida cotidiana trato de limitar el uso de herramientas y plataformas que capturan grandes volúmenes de información, pero la omnipresencia de estos servicios lo dificulta.
La lectura deja como aportación valiosa la claridad sobre cómo somos el combustible gratuito de las tecnológicas.
Sin embargo, el tema abre un panorama sombrío frente al futuro: la pregunta y el reto mayor para mí es que, después de 8 años desde la publicación del artículo, no veo una mejora sino un severo agravamiento del problema.
Frente al poder de cálculo y la voracidad de las grandes empresas de IA y Big Data modernas, no logro ver cómo alcanzaremos efectivamente ese empoderamiento ciudadano que nos permita reclamar nuestra soberanía digital.
